\part{更深一层的思考：与这个世界相处}

\section{思考：一个更舒服的看待世界的角度}

方法讲完了。这一章不讲操作，讲一个更舒服的看待世界的角度——把前面那些散落的认知，拼成一张完整的世界观。

\subsection{为什么学校会惩罚 ADHD}

先说一个很多人没想明白的问题：为什么 ADHD 在学校里总是吃亏？

答案不是「学校更难」，而是\term{学校默认深度优先算法}——一节课四十五分钟、一章接一章、线性积累、一次没跟上就永远跟不上。这套运行模式恰好踩在 ADHD 算法最贵的那块成本上。上课听不懂，浪费时间，产生挫败感；搜寻范围过大，缺乏粗糙分类，脑内一团浆糊；突发兴趣转不过去，想法无法深入发展。

\begin{insight}
注意这里的关键：\term{这不是你的错，是环境不匹配}。你在课堂里「失败」，和你拿着 GPU 去跑单线程渲染失败，性质是一样的——工具没错，是任务和工具不匹配。把「我上课听不进去」翻译成「我在一个不适合我的运行环境里」，你就不会再用「我笨」「我没救」来惩罚自己。
\end{insight}

\subsection{「轴的程度」与「适不适合」：维度不匹配不是病}

原视频讲过一个小案例，非常能说明问题（BV13HoGYiEgF）：一个能考上北大光华的学生，高等数学考三四十分，怎么教都不会。你很难说他是笨——能上光华的人，怎么会连高数都学不会？更合理的解释是：\term{某个维度不匹配}。

原视频还有个更狠的类比：人的「轴的程度」各不相同——轴到九十九分的人不适合做数学，轴到四五十分的人也不适合，八十分左右最适合。那是不是该给「高轴」的人判一个「高轴综合症」、给「低轴」的人判一个「低轴综合症」？这当然是搞笑——\term{谁规定你生下来就得做数学了}？如果按「适不适合做数学」来给人分类，多分几个维度，每个人都有一堆病。

所以对「我是不是有缺陷」这个问题，一个更舒服的回答是：\term{缺陷与否，取决于它在你生活里造成了多大困扰}。特征 + 损害 = 诊断；没有损害，就只有特征。方法解决得了的困扰，就不是「病」，是「待适配」。

\subsection{疾病观的延伸：方法优先，药物兜底}

顺着这个逻辑，原视频对「ADHD 算不算病」的态度很明确（BV13HoGYiEgF）：如果你靠布局神经元、连接整理那一套手法，依然觉得费劲、难以正常工作，那你确实需要药物，不必硬扛；但如果方法能解决，那就没必要把它当病治。

这里的核心立场是\term{实用主义}：不问「我算不算病」，只问「什么让我运转得最好」。方法能解决的用方法，方法解决不了的，药物和正规医疗是合理且正当的兜底——\term{用工具解决问题，不丢人}。这两条路不冲突，更不是道德高下之分。

\subsection{以广促深：广度算法在长周期里能赢}

ADHD 的广度算法，在短周期、单任务的视角下是「缺陷」；但放到长周期、多任务的视角下，它有独特的优势。

\term{深度运算虽然贵，但「以广促深」是成立的}：你了解并领悟很多不同的视角，切换着看问题，最后能触达的深度，不一定比线性深挖的人浅。很多重要发现恰恰来自「把 A 领域的方法搬进 B 领域」——而 ADHD 的广度搜索，天生擅长制造这种跨界连接。

原视频甚至有个大胆的猜想：\term{欧拉可能是 ADHD}。欧拉晚年失明了，依然能不停口述写作、产出海量论文——因为他的写作是「默认模式」，靠被动注意力驱动，根本不需要控制。\term{这是广度算法被逼到极限之后的样子：不需要开关，一直在产出}。

\begin{insight}
所以「ADHD 是不是缺陷」，最后的答案是这样的：它是一套在当下环境里处处吃亏、但在长周期里藏着独特优势的算法。前期的关卡难打，后期的优势真实存在。你不需要为「广度」道歉，你需要的是把它用对地方。
\end{insight}

\subsection{关于「正常人」：他们也被极大高估了}

整份讲义里，我们反复提到「普通人」「正常人」。最后要提醒一句：\term{他们也被极大高估了}。

你羡慕的那个人，可能只是恰好活在适合他的环境里；你觉得自己「连普通人都不如」，可能只是你拿自己最弱的运行模式，去比人家最强的运行模式。多分几个维度，每个人都有「病」，也每个人都有「天赋」。

原视频用一句很朴素的话收过这个主题：\term{不要老想着别人能做到，你也一定能做到}——这句话反过来也成立：\term{别人做到了，也不代表你必须用同一条路做到}。你可以最后做出好的作品，但过程不需要和别人一样。孙子那句话用在这里很合适：\term{能屈能伸，方为丈夫}。

\subsection{给 ADHD 群体的一句话}

把这一切收束成一句给同路人的话，原视频的说法是：

\term{先混过去，再放下，再重新自发地产生兴趣。相信顺应冲动，也能做出很多事情。}

你不是需要被治愈，你是需要被读懂。这份讲义，就是读懂自己的开始。
